% arXiv-style skeleton. Mirrors ../WRITEUP.md; numbers come only from the notebook.
% Build: latexmk -pdf main.tex   (not built in CI; no TeX toolchain is assumed)
\documentclass[11pt]{article}
\usepackage[margin=1in]{geometry}
\usepackage{amsmath,booktabs,hyperref,natbib,xcolor}
\newcommand{\TODO}[1]{\textbf{\textcolor{red}{TODO: #1}}}

\title{Memory Anchored to the Code Graph Helps Small Coding Agents?}
\author{TODO authors}
\date{}

\begin{document}
\maketitle

\begin{abstract}
Experience memory lets an agent reuse lessons from earlier attempts, but it is usually retrieved
by text similarity. We anchor each lesson to a node of the repository's AST call graph and
retrieve it through the call-graph neighbourhood and precomputed node embeddings, which lets a
lesson written on one repository reach similar code in an unseen one. Each Gemma~4 model (E4B,
12B, 31B) distils lessons only from its own attempts. With the lesson pool, top-$k$ and context
budget held fixed, we compare no memory, flat text-similarity memory and graph-anchored memory,
leave-one-repo-out on the 129 development tasks of the Gemma~4 Developer Agent harness, under a
decision rule preregistered before any run. \TODO{verdict, pooled difference and 95\% CI}
\end{abstract}

\section{Introduction}
Reflexion~\citep{shinn2023reflexion}, ExpeL~\citep{zhao2023expel}, Agent Workflow
Memory~\citep{wang2024awm} and ReasoningBank~\citep{ouyang2025reasoningbank} keep lessons from an
agent's own trajectories and retrieve them by text or task similarity. For code, the relevant key
is the code a task touches. \TODO{motivation paragraph; contributions list as in WRITEUP.md}

\section{Method}
\subsection{Lessons}
A lesson is $\ell = (v, r, t, o, \tau, m)$: anchor node $v$, repository $r$, text $t$
($\le 400$ characters), outcome $o \in \{\text{success}, \text{failure}\}$, source task $\tau$
and source model $m$. The model $m$ that made the attempt writes $\ell$; $v$ must be a node the
attempt touched.

\subsection{Graph-anchored retrieval}
Let $S$ be the task's seed nodes, $d(s, v)$ the call-graph hop distance and $e_v$ the harness's
node embedding. A lesson anchored at $v$ scores
\begin{equation}
  \mathrm{score}(\ell) = \max_{s \in S}
  \begin{cases}
    \gamma^{d(s,v)} & d(s,v) \le H,\\
    \lambda \cos(e_s, e_v) & d(s,v) > H \text{ and } \cos(e_s, e_v) \ge \theta,
  \end{cases}
\end{equation}
with $\gamma = 0.6$, $H = 2$, $\lambda = 0.8$ and $\theta = 0.5$, all fixed before any run.
The top $k = 5$ lessons are kept under a 1{,}600-character budget, with at most two per anchor.
The flat arm uses BM25 over the lesson text with the same $k$, budget and rendering.

\section{Experiments}
\subsection{Protocol}
\TODO{models, LORO folds, hidden test set, preregistered rule (paired bootstrap clustered by task,
exact sign test, Holm)}

\subsection{Results}
\begin{table}[h]
\centering
\caption{Success rate, development tasks, leave-one-repo-out, harness scorer.}
\begin{tabular}{lccc}
\toprule
Model & none & flat & graph \\
\midrule
gemma-4-e4b-it-qat-w4a16 & \TODO{} & \TODO{} & \TODO{} \\
gemma-4-12b-it-qat-w4a16 & \TODO{} & \TODO{} & \TODO{} \\
gemma-4-31b-it-qat-w4a16 & \TODO{} & \TODO{} & \TODO{} \\
\bottomrule
\end{tabular}
\end{table}

\TODO{Table: H1 and H2 contrasts; Table: hidden-test leaderboard scores; Table: cost}

\subsection{Threats to validity}
\TODO{as in WRITEUP.md section 3.5}

\section{Related work}
Code graphs are used to retrieve code by RepoGraph~\citep{ouyang2024repograph},
CodexGraph~\citep{liu2024codexgraph}, LocAgent~\citep{chen2025locagent} and
\citet{bhola2026codeisntmemory}; coding agents follow SWE-agent~\citep{yang2024sweagent}.
EA-Graph~\citep{hsu2026eagraph} anchors verification claims to artifacts and is the closest prior
work; \citet{wang2025repomemory} mine commit history for memory, and SWE-Exp~\citep{chen2025sweexp}
keeps an experience bank of repair attempts. General memory systems include
MemGPT~\citep{packer2023memgpt}, Mem0~\citep{chhikara2025mem0}, Zep~\citep{rasmussen2025zep} and
Hindsight~\citep{latimer2025hindsight}, evaluated on benchmarks such as
LongMemEval~\citep{wu2024longmemeval}; VibeMemBench~\citep{fan2026vibemembench} finds most memory
systems do not beat no memory for coding agents.

\section{Conclusion}
\TODO{verdict in the preregistered words, cost, main limitation}

\bibliographystyle{plainnat}
\bibliography{refs}
\end{document}
